\documentclass[10pt,a4paper]{article}

\usepackage[a4paper,margin=20mm]{geometry}
\usepackage[T1]{fontenc}
\usepackage[utf8]{inputenc}
\usepackage{lmodern}

\usepackage{graphicx}
\usepackage{booktabs}
\usepackage{tabularx}
\usepackage{array}
\usepackage{longtable}
\usepackage{float}
\usepackage{enumitem}
\usepackage{fancyhdr}
\usepackage{titlesec}
\usepackage{hyperref}
\usepackage{xcolor}
\usepackage{listings}
\usepackage{amsmath}
\usepackage{multirow}

% ============================================================
% HOW TO USE THIS TEMPLATE
% ============================================================
% 1. Fill in the \IPName, \IPVersion, \IPAuthor, \IPDate commands below.
% 2. Replace every <ANGLE-BRACKET> placeholder with real content.
% 3. Duplicate table rows / itemize entries as needed for your IP's
%    actual ports, parameters, tests, etc. Delete unused example rows.
% 4. Sections 6 (Verification) and 7 (Synthesis) are left with "--"
%    placeholders on purpose — fill in once results are available,
%    or delete rows that don't apply to your IP.
% 5. Save this file as <ip_name>_doc.tex and rename the block diagram
%    image accordingly.
% ============================================================

% ============================================================
% Document Information -- EDIT THESE FOR EACH NEW IP
% ============================================================
\newcommand{\IPName}{<IP Name>}
\newcommand{\IPVersion}{v0.1}
\newcommand{\IPAuthor}{<Author Name>}
\newcommand{\IPDate}{<Month Year>}

\title{
    \vspace{-20mm}
    \textbf{\IPName}\\[2mm]
    \large IP Design Documentation
}

\author{\IPAuthor}
\date{\IPDate}

% ============================================================
% Page Style
% ============================================================
\pagestyle{fancy}
\fancyhf{}

\lhead{\textbf{\IPName}}
\rhead{\IPVersion}
\cfoot{\thepage}

\setlength{\headheight}{14pt}

% ============================================================
% Section Style
% ============================================================
\titleformat{\section}
    {\large\bfseries}
    {\thesection.}
    {0.5em}
    {}

\titleformat{\subsection}
    {\normalsize\bfseries}
    {\thesubsection}
    {0.5em}
    {}

% ============================================================
% Hyperlink
% ============================================================
\hypersetup{
    colorlinks=true,
    linkcolor=black,
    urlcolor=blue,
    citecolor=black
}

% ============================================================
% Listings
% ============================================================
\lstset{
    basicstyle=\ttfamily\small,
    frame=single,
    breaklines=true,
    columns=fullflexible
}

% ============================================================
% Document
% ============================================================
\begin{document}

% ============================================================
% Cover Page
% ============================================================

\begin{titlepage}
\centering
\vspace*{\fill}
{\Huge\bfseries \IPName\par}\vspace{8mm}
{\Large IP Design Documentation\par}\vspace{15mm}
{\large Version \IPVersion\par}\vspace{3mm}
{\large\IPAuthor\par}\vspace{3mm}
{\large \IPDate\par}\vspace*{\fill}

\end{titlepage}

% ============================================================
\section{Revision}
% ============================================================

\begin{table}[H]
\centering
\begin{tabularx}{\textwidth}{c c l X}
\toprule
\textbf{Version} & \textbf{Date} & \textbf{Author} & \textbf{Description} \\
\midrule
v0.1 & <Month Year> & <Author Name> & Initial release \\
% Add a new row per revision, e.g.:
% v0.2 & <Month Year> & <Author Name> & <What changed> \\
\bottomrule
\end{tabularx}
\end{table}

% ============================================================
\section{Overview}
% ============================================================

\subsection{Purpose}

% One short paragraph: what does this IP do, and in what context
% is it used (e.g. clock domain, bus protocol, system role)?
<Describe the purpose of the IP in 2-4 sentences.>

\subsection{Features}

\begin{itemize}[noitemsep]
    \item <Feature 1>
    \item <Feature 2>
    \item <Feature 3>
    % Add/remove bullets as needed
\end{itemize}


% ============================================================
\section{Architecture}
% ============================================================

\subsection{Block Diagram}

\begin{figure}[H]
    \centering
    \includegraphics[width=\textwidth]{<block_diagram_filename>.png}
    \caption{<IP Name> block diagram}
    \label{fig:block_diagram}
\end{figure}

\subsection{IO Ports}

% Group ports logically (Clock/Reset, Control, Data, Status, etc.)
% and add/remove rows as needed. \multirow{N}{*}{Group} spans N rows.

\begin{table}[H]
\centering
\renewcommand{\arraystretch}{1.2}
\begin{tabularx}{\textwidth}{|c|c|c|c|X|}
\hline
\textbf{Group} &
\textbf{Port} &
\textbf{Direction} &
\textbf{Width} &
\multicolumn{1}{c|}{\textbf{Description}} \\
\hline

\multirow{2}{*}{Clock/Reset}
    & clk\_i  & Input  & 1  & Clock signal \\
    & rst\_i  & Input  & 1  & Asynchronous reset \\
\hline

\multirow{2}{*}{<Group Name>}
    & <port\_name\_i>  & Input  & <WIDTH>  & <Description> \\
    & <port\_name\_o>  & Output & <WIDTH>  & <Description> \\
\hline

% Duplicate the \multirow block above for each additional port group

\end{tabularx}
\caption{<IP Name> interface}
\label{tab:interface}
\end{table}

\subsection{Config Parameters}

\begin{table}[H]
\centering
\renewcommand{\arraystretch}{1.2}

\begin{tabularx}{\textwidth}{|c|c|c|X|}
\hline
\textbf{Parameter} &
\textbf{Default} &
\textbf{Range} &
\multicolumn{1}{c|}{\textbf{Description}} \\
\hline

<PARAM\_NAME> & <default> & <range> & <Description> \\
\hline
% Add one row per parameter

\end{tabularx}

\caption{Configurable parameters}
\label{tab:parameters}
\end{table}


% ============================================================
\section{Functional Description}
% ============================================================
 
% Describe every functional capability of the IP and how to operate
% it, referencing exact port names. There is no fixed subsection
% breakdown here -- write as many or as few subsections as the IP
% actually needs (e.g. some IPs need just one operation description,
% others need several: normal operation, configuration modes, error
% handling, reset, status, etc.). Suggested topics to cover somewhere
% in this section, as applicable:
%   - Overall operating behavior end to end
%   - Each distinct operation/mode and the conditions that trigger it
%   - Reset behavior (assertion/deassertion, resulting state)
%   - Any special behaviors: flush, error handling, back-pressure,
%     arbitration, power modes, etc.
%   - Status/flag outputs and what they indicate
% Delete the example subsection below and replace with real content.
 
\subsection{<Descriptive subsection title as needed>}
 
<Describe this piece of functional behavior/operation guide,
referencing exact port names. Add or remove subsections freely
to fully cover the IP's functional capability.>


% ============================================================
\section{Verification Test}
% ============================================================
 
% Simple sanity checks confirming the IP works as intended --
% not full DV coverage. List the test cases that were run
% and passed.
\begin{table}[H]
\centering
\begin{tabularx}{\textwidth}{l c X}
\toprule
\textbf{Test} & \textbf{Status} & \textbf{Description} \\
\midrule
<Test name> & PASS & <What the test checks> \\
% Add one row per test case
\bottomrule
\end{tabularx}
\end{table}

% ============================================================
\section{Synthesis Report}
% ============================================================
 
% Synthesized using the Nangate45 library. Fill in the results below.
\begin{table}[H]
\centering
\begin{tabularx}{\textwidth}{l X}
\toprule
\textbf{Metric} & \textbf{Result} \\
\midrule
Library     & Nangate45 \\
Frequency   & -- \\
Cell Count  & -- \\
Cell Area   & -- \\
WNS         & -- \\
\bottomrule
\end{tabularx}
\end{table}
 
% Only include an SDC subsection if this IP needs special timing
% constraints (e.g. multicycle paths, false paths, generated clocks).
% Otherwise, delete the subsection below.
\subsection{SDC}
 
\begin{lstlisting}
<constraint command, e.g. set_multicycle_path 2 -from ... -to ...>
\end{lstlisting}
 
 
% ============================================================
\section{Note}
% ============================================================
 
% Optional:
% - Known limitations
% - Integration notes
% - Design assumptions
% - Known issues
<Add any notes here, or delete this section if not needed.>

\end{document}
